\subsubsection{\texorpdfstring{$\Xi_\zeta$}{Xi\_zeta} 的单点喷射到 Verblunsky--CMV 反射谱：Toeplitz 证书的标量化与 LMI 提升}\label{subsubsec:xi-verblunsky-cmv-from-local-jet}
本段把“视界对数导数的正性 \(\Leftrightarrow\) Toeplitz--PSD 层级”（附录定理 \ref{thm:app-horizon-toeplitz-lmi}）进一步标量化为一列反射系数（Verblunsky/Schur 参数），并强调其\textbf{单点局部闭包}：在每一有限级别上，反射系数与相应 CMV（五对角）对象都可由 \(\Xi_\zeta\) 在 \(s=-\tfrac12\) 的有限阶对数导数数据 \(\{\ell_k\}\) 显式生成（命题 \ref{prop:app-horizon-coeff-local}）。
由此得到一个可审计的“逆谱接口”：
\[
\boxed{\;
(\ell_k)_{k\ge 1}\ \Longrightarrow\ (c_n)_{n\ge 0}\ \Longrightarrow\ (\alpha_n)_{n\ge 0}
\ \Longrightarrow\ 
\text{CMV（酉）反射谱对象}\;}
\]
并使 Toeplitz 层级的失败定位从“矩阵负特征值存在”进一步压缩为“最小反射系数越界”的单标量事件。

\paragraph{从 Carath\'eodory 到 Schur：规范反射谱的定义}
令 \(\mathscr{C}_{\zeta}\) 为视界 Carath\'eodory 对数导数函数（附录式 \eqref{eq:app-horizon-carath-def}），并写其幂级数
\[
\mathscr{C}_{\zeta}(w)=c_0+2\sum_{n\ge 1}c_n w^n,
\qquad c_{-n}:=\overline{c_n}.
\]
在 RH 成立时，\(\mathscr{C}_\zeta\) 属于 Carath\'eodory 类且 \(c_0=\Re\mathscr{C}_\zeta(0)=\mu_\zeta(\TT)>0\)（推论 \ref{cor:app-horizon-carath-rh}；并见附录定理 \ref{thm:app-horizon-spectral-measure-atomic} 的测度解释）。
此外由自对偶对称可取实结构 \(c_{-n}=c_n\in\RR\)（推论 \ref{cor:app-horizon-ramanujan-selfdual-measure}），从而 Toeplitz 主子矩阵 \(\mathbf{T}_N=(c_{j-k})_{0\le j,k\le N}\) 为实对称矩阵。
\par
为把 Toeplitz--PSD 层级压缩为标量递归，令
\[
F(w):=\frac{\mathscr{C}_{\zeta}(w)}{c_0},
\qquad F(0)=1.
\]
并定义相应的 Schur 函数
\begin{equation}\label{eq:xi-schur-from-carath}
\boxed{\;
S(w):=\frac{1}{w}\frac{F(w)-1}{F(w)+1}\; }.
\end{equation}
当 \(F\) 为 Carath\'eodory 时，\(S\) 为 Schur 类函数（\(|S(w)|\le 1\) 于 \(\mathbb{D}\)；参见 \cite{Simon2005OPUC1}）。
因此存在唯一 Verblunsky 系列 \((\alpha_n)_{n\ge 0}\subset\overline{\mathbb{D}}\) 使 Schur 递归成立：
\[
S_0=S,\qquad \alpha_n:=S_n(0),\qquad
S_{n+1}(w):=\frac{1}{w}\frac{S_n(w)-\alpha_n}{1-\overline{\alpha_n}\,S_n(w)}.
\]
在 RH 下，由 \(\mathbf{T}_N\succeq 0\) 对一切 \(N\)（定理 \ref{thm:app-horizon-toeplitz-lmi}）推出 \(|\alpha_n|\le 1\)（见下述定理 \ref{thm:xi-toeplitz-det-verblunsky}），从而 \((\alpha_n)\) 构成体系内的规范“反射谱”。

\begin{theorem}[\texorpdfstring{$\Xi_\zeta$}{Xi\_zeta} 的规范反射谱：\texorpdfstring{$\ell$}{ell}\,--\,\texorpdfstring{$c$}{c}\,--\,\texorpdfstring{$\alpha$}{alpha} 的有限阶闭式]\label{thm:xi-verblunsky-from-local-jet}
\leanverified{paper\_xi\_verblunsky\_from\_local\_jet}
在上述记号下，设 \(c_0>0\) 并令 \(S\) 由 \eqref{eq:xi-schur-from-carath} 定义。
则：
\begin{enumerate}
  \item \(\alpha_0=S(0)\) 与 \(\alpha_1=S_1(0)\) 由 \((c_0,c_1,c_2)\) 给出严格闭式
  \[
  \boxed{\;
  \alpha_0=\frac{c_1}{c_0},\qquad
  \alpha_1=\frac{c_0c_2-c_1^2}{c_0^2-c_1^2}\;}
  \]
  （在实结构下 \(c_1\in\RR\)；一般复结构时分母应写为 \(c_0^2-|c_1|^2\)）。
  \item 若进一步用单点喷射
  \(
  \ell_k=\left.\frac{\dd^k}{\dd s^k}\log\Xi_\zeta(s)\right|_{s=-1/2}
  \)
  表示局部数据，则由命题 \ref{prop:app-horizon-coeff-local} 的 \(c_0=-\ell_1,c_1=-\ell_2,c_2=\ell_2-\ell_3\) 得
  \[
  \boxed{\;
  \alpha_0=\frac{c_1}{c_0}=\frac{\ell_2}{\ell_1},\qquad
  \alpha_1=\frac{-\ell_1\ell_2+\ell_1\ell_3-\ell_2^2}{\ell_1^2-\ell_2^2}\;}
  \]
  \item 更一般地，对每个 \(n\ge 0\)，\(\alpha_n\) 是 \((c_0,\dots,c_{n+1})\) 的有理函数，从而亦是 \((\ell_1,\dots,\ell_{n+2})\) 的有理函数；因此反射谱在任一有限级别上由单点局部数据完全生成。
\end{enumerate}
\end{theorem}

\begin{proof}
由 \(F(0)=1\) 可知 \(\frac{F(w)-1}{F(w)+1}\) 在 \(w=0\) 处一阶消失，故 \eqref{eq:xi-schur-from-carath} 定义了解析的 \(S\)。
将 \(F(w)=1+2(c_1/c_0)w+2(c_2/c_0)w^2+O(w^3)\) 代入并作幂级数除法，得到
\[
S(w)=\frac{c_1}{c_0}+\left(\frac{c_0c_2-c_1^2}{c_0^2-c_1^2}\right)w+O(w^2),
\]
从而 \(\alpha_0=S(0)=c_1/c_0\)。
再由 Schur 递归
\(
S_1(w)=\frac{1}{w}\frac{S(w)-\alpha_0}{1-\overline{\alpha_0}\,S(w)}
\)
取 \(w\to 0\) 的极限，即得 \(\alpha_1=S_1(0)\) 的闭式；实结构下 \(\overline{\alpha_0}=\alpha_0\) 给出所示形式。
第二条由命题 \ref{prop:app-horizon-coeff-local} 代换得到。
第三条由 Schur 算法：每一步 \(S_{n+1}\) 由 \(S_n\) 通过 Möbius 变换与除以 \(w\) 得到，故其在 \(w=0\) 的值 \(\alpha_{n+1}=S_{n+1}(0)\) 仅依赖于 \(S_n\) 的有限阶 Taylor 系数；
而 \(S\) 的 \(n\) 阶 Taylor 系数仅依赖于 \(F\) 的 \(n{+}1\) 阶系数，进而仅依赖于 \((c_0,\dots,c_{n+1})\)（以及实结构下的共轭对称性）。
最后由命题 \ref{prop:app-horizon-coeff-local}，每个 \(c_k\) 由有限个 \(\ell\)-喷射数据有理生成，从而 \(\alpha_n\) 亦然。证毕。
\end{proof}

\begin{corollary}[二阶喷射的有限阶否证见证：\texorpdfstring{$(\alpha_0,\alpha_1)$}{(alpha0,alpha1)} 的越界判据]\label{cor:xi-horizon-order2-reflection-falsifier}
\leanpartial{paper\_xi\_horizon\_order2\_reflection\_falsifier}{the requested Lean signature is false in the current repository state: with \(c_0=c_1=1\) and \(c_2=5\), the repository definition gives \(\alpha_0=1\) and \(\alpha_1=0\), so the hypothesis holds but the displayed conclusion becomes \(4 \le 0\).}
沿用定理 \ref{thm:xi-verblunsky-from-local-jet} 的记号，并假设 \(c_0>0\)。
若 Toeplitz--PSD 层级对该喷射成立（从而处于 RH 相的 Carath\'eodory/Schur 口径），则必须有
\[
\boxed{\;
|c_1|\le c_0,\qquad
|c_0c_2-c_1^2|\le c_0^2-|c_1|^2.\;}
\]
因此一旦出现任一违背
\[
|c_1|>c_0
\quad\text{或}\quad
|c_0c_2-c_1^2|>c_0^2-|c_1|^2,
\]
则有限阶 Toeplitz 证书必在 \(N\le 2\) 的层级上失败，从而给出可离线复核的有限否证。
\end{corollary}
\begin{proof}
Toeplitz--PSD 蕴含 \(|\alpha_n|\le 1\)（定理 \ref{thm:xi-toeplitz-det-verblunsky}），特别地 \(|\alpha_0|\le 1\)。
由定理 \ref{thm:xi-verblunsky-from-local-jet} 的闭式
\(
\alpha_0=c_1/c_0
\)
得到 \(|c_1|\le c_0\)。
当 \(|c_1|<c_0\) 时分母 \(c_0^2-|c_1|^2>0\)，并可由同一定理的闭式
\(
\alpha_1=(c_0c_2-c_1^2)/(c_0^2-|c_1|^2)
\)
（一般复结构时的分母形式），即可得到所示两条不等式。逆否命题给出有限阶失败见证。证毕。
\end{proof}

\begin{corollary}[Toeplitz--PSD 的喷射标量化：失败可消元为有限个显式代数不等式]\label{cor:xi-toeplitz-psd-jet-semialgebraic}
\leanverified{paper\_xi\_toeplitz\_psd\_jet\_semialgebraic}
沿用定理 \ref{thm:xi-verblunsky-from-local-jet} 的记号，固定 \(N\ge 1\)，并在非退化口径下假设
\(\det\mathbf{T}_m>0\) 对所有 \(0\le m\le N\) 成立。
则 \(\mathbf{T}_N\succeq 0\) 等价于 Verblunsky 前缀满足
\[
|\alpha_k|\le 1\qquad(0\le k\le N-1).
\]
并且每个 \(\alpha_k\) 都是有限喷射 \((\ell_1,\dots,\ell_{k+2})\) 的有理函数。
因此，上述不等式组在逐项写成 \(|\alpha_k|^2\le 1\) 并清分母后，等价于关于有限喷射 \((\ell_1,\dots,\ell_{N+2})\) 的有限个显式多项式不等式系统。
从而“某一级 Toeplitz--PSD 失败”可被转写为一个有限维半代数可行性问题，否证器搜索的计算对象可由大矩阵谱问题压缩为低维标量不等式系统。
\end{corollary}
\begin{proof}
Schur 算法的有限步可行性判据给出：在 \(\det\mathbf{T}_m>0\) 的非退化口径下，\(\mathbf{T}_N\succeq 0\) 当且仅当全部前缀反射系数满足 \(|\alpha_k|\le 1\)（\(k<N\)）。
而 \(\alpha_k\) 对喷射 \((\ell_1,\dots,\ell_{k+2})\) 的有理依赖由定理 \ref{thm:xi-verblunsky-from-local-jet} 给出。
将每个约束写为 \(|\alpha_k|^2\le 1\) 并清分母，即得到关于 \((\ell_1,\dots,\ell_{N+2})\) 的有限个多项式不等式。证毕。
\end{proof}

\paragraph{Toeplitz--PSD 层级的反射乘积分解与最小失败步定位}
反射系数不仅给出“正性”的标量化参数，而且给出 Toeplitz 行列式的乘积分解，从而把有限维否证器定位到最小越界步。

\begin{theorem}[Toeplitz 行列式的反射乘积分解与最小失败索引]\label{thm:xi-toeplitz-det-verblunsky}
\leanverified{paper\_xi\_toeplitz\_det\_verblunsky}
沿用上述记号，并令 \(\mathbf{T}_N=(c_{j-k})_{0\le j,k\le N}\)。
若 \(\mathbf{T}_N\) 正定，则存在唯一 \(\alpha_0,\dots,\alpha_{N-1}\in\mathbb{D}\)（Verblunsky 系列的前 \(N\) 项）使得对所有 \(1\le m\le N\) 有
\begin{equation}\label{eq:xi-toeplitz-det-product}
\boxed{\;
\det \mathbf{T}_m
=c_0^{m+1}\prod_{k=0}^{m-1}\bigl(1-|\alpha_k|^2\bigr)^{\,m-k}\; }.
\end{equation}
因此 \(\det \mathbf{T}_m>0\) 的充要条件为 \(|\alpha_k|<1\) 对所有 \(0\le k\le m-1\) 成立。
反之，若存在最小 \(m_\star\ge 1\) 使 \(\mathbf{T}_{m_\star}\not\succeq 0\)，则 Schur 递归在 \(m_\star-1\) 步之前严格可行，并且在最小失败步上必出现
\[
|\alpha_{m_\star-1}|\ge 1,
\qquad\text{从而}\qquad
\det \mathbf{T}_{m_\star}\le 0.
\]
\end{theorem}

\begin{proof}
式 \eqref{eq:xi-toeplitz-det-product} 是 OPUC/Schur 算法与 Levinson--Durbin 因子分解的标准推论：\(\mathbf{T}_m\) 的 Cholesky 分解可由反射系数递归构造，且每一步引入因子 \(1-|\alpha_k|^2\)（参见 \cite[Ch.~1--2]{Simon2005OPUC1} 或截断矩问题的 Carath\'eodory--Fej\'er 理论）。
由此得到 \(\det \mathbf{T}_m>0 \Longleftrightarrow |\alpha_k|<1\)（\(k<m\)）。
反向定位：若 \(\mathbf{T}_{m_\star}\not\succeq 0\) 且 \(m_\star\) 最小，则对所有 \(m<m_\star\) 有 \(\mathbf{T}_m\succeq 0\) 且（在 \(c_0>0\) 下）\(\mathbf{T}_m\succ 0\)，从而 Schur 递归在前 \(m_\star-1\) 步可行并产生 \(\alpha_0,\dots,\alpha_{m_\star-2}\in\mathbb{D}\)。
若仍有 \(|\alpha_{m_\star-1}|<1\)，则由 \eqref{eq:xi-toeplitz-det-product} 得 \(\det\mathbf{T}_{m_\star}>0\)，与 \(\mathbf{T}_{m_\star}\not\succeq 0\) 矛盾；故必须 \(|\alpha_{m_\star-1}|\ge 1\)，并使 \(\det\mathbf{T}_{m_\star}\le 0\)。证毕。
\end{proof}

\begin{corollary}[实现维数阈值强制的反射越界层级上界]\label{cor:xi-toeplitz-reflection-overrun-index-bound}
\leanverified{paper\_xi\_toeplitz\_reflection\_overrun\_index\_bound}
在定理 \ref{thm:xi-toeplitz-psd-coherence-horizon-threshold} 的“实现维数 \(\le D\)”有界类设定下，设 \(N_0\le 2D\) 为其所给 Toeplitz--PSD 视界阈值：
\[
\mathbf{T}_N\succeq 0\ (\forall N\ge 0)
\quad\Longleftrightarrow\quad
\mathbf{T}_N\succeq 0\ (0\le N\le N_0).
\]
若对某个参数实例 Toeplitz--PSD 失败，并令 \(m_\star\ge 1\) 为最小失败层级
\(
\mathbf{T}_{m_\star}\not\succeq 0
\)，
则必有
\[
\boxed{\;
m_\star\ \le\ N_0\ \le\ 2D\;}
\qquad\text{且}\qquad
|\alpha_{m_\star-1}|\ge 1.
\]
\end{corollary}
\begin{proof}
Toeplitz--PSD 失败意味着存在某个 \(N\ge 0\) 使 \(\mathbf{T}_N\not\succeq 0\)。
由定理 \ref{thm:xi-toeplitz-psd-coherence-horizon-threshold} 的视界等价性，可取这样的 \(N\) 满足 \(N\le N_0\)，从而最小失败层级 \(m_\star\le N\le N_0\le 2D\)。
最后由定理 \ref{thm:xi-toeplitz-det-verblunsky} 的最小失败步定位得 \(|\alpha_{m_\star-1}|\ge 1\)。证毕。
\end{proof}

\begin{corollary}[“视界主子式”一致性：Toeplitz 行列式序列的对数凹性与反射参数界]\label{cor:xi-toeplitz-principal-minor-logconcavity}
\leanverified{paper\_xi\_toeplitz\_principal\_minor\_logconcavity}
沿用定理 \ref{thm:xi-toeplitz-det-verblunsky} 的记号，令 \(\Delta_N:=\det\mathbf{T}_N\)（\(N\ge 0\)）。
当 \(\mathbf{T}_N\succ 0\) 对所有 \(N\) 成立时，行列式序列 \((\Delta_N)_{N\ge 0}\) 满足显式比值恒等式
\[
\boxed{\;
\frac{\Delta_{N+1}\Delta_{N-1}}{\Delta_N^{2}}=1-|\alpha_N|^2\qquad(N\ge 1),\;}
\]
从而 \((\Delta_N)\) 为对数凹序列：
\[
\boxed{\;
\Delta_N^2\ \ge\ \Delta_{N-1}\Delta_{N+1}\qquad(N\ge 1).\;}
\]
特别地，\(\Delta_{N+1}\Delta_{N-1}/\Delta_N^2\in(0,1]\) 等价于 \(|\alpha_N|\le 1\)；
因此在非退化情形（\(\Delta_N>0\)）下，\(|\alpha_N|\le 1\)（所有 \(N\)）可被视为 Toeplitz--PSD 层级的一组纯标量一致性检查：
它把“全级正性”压缩为“主子式的对数凹性（跨级一致性）+ 行列式正性（级内可行性）”两类单调约束。
\end{corollary}
\begin{proof}[证明要点]
由 \eqref{eq:xi-toeplitz-det-product} 直接计算
\[
\frac{\Delta_{N+1}}{\Delta_N}
=c_0\prod_{k=0}^{N}\bigl(1-|\alpha_k|^2\bigr),
\qquad
\frac{\Delta_{N}}{\Delta_{N-1}}
=c_0\prod_{k=0}^{N-1}\bigl(1-|\alpha_k|^2\bigr),
\]
两式相除即得 \(\Delta_{N+1}\Delta_{N-1}/\Delta_N^2=1-|\alpha_N|^2\)。
对数凹性由 \(1-|\alpha_N|^2\le 1\) 立得。证毕。
\end{proof}

\begin{theorem}[Toeplitz 主子式的离散曲率公式与 Verblunsky 模长的完全恢复]\label{thm:xi-toeplitz-principal-minor-discrete-curvature-recovery}
\leanverified{paper\_xi\_toeplitz\_principal\_minor\_discrete\_curvature\_recovery}
沿用定理 \ref{thm:xi-toeplitz-det-verblunsky} 的记号，并记
\[
\Delta_m:=\det\mathbf{T}_m\qquad(m\ge 0).
\]
在非退化口径 \(\Delta_m>0\)（\(m\ge 0\)）下，对每个 \(m\ge 1\) 都有严格恒等式
\begin{equation}\label{eq:xi-toeplitz-principal-minor-discrete-curvature}
\boxed{\;
1-|\alpha_m|^2
=
\frac{\Delta_{m+1}\Delta_{m-1}}{\Delta_m^2}\; }.
\end{equation}
等价地，主子式序列的离散曲率可写为
\[
\mathcal K_m:=-\log\!\left(\frac{\Delta_{m+1}\Delta_{m-1}}{\Delta_m^2}\right)
=-\log(1-|\alpha_m|^2)
\qquad(m\ge 1).
\]
从而成立：
\begin{enumerate}
  \item 对每个 \(m\ge 1\)，都有
  \[
  \Delta_m^2\ge \Delta_{m-1}\Delta_{m+1},
  \]
  且严格不等式成立当且仅当 \(|\alpha_m|>0\)；
  \item 全体模长序列 \((|\alpha_m|)_{m\ge 0}\) 仅由主子式序列 \((\Delta_m)_{m\ge 0}\) 唯一决定，其中
  \[
  1-|\alpha_0|^2=\frac{\Delta_1}{\Delta_0^2},
  \qquad
  1-|\alpha_m|^2=\frac{\Delta_{m+1}\Delta_{m-1}}{\Delta_m^2}\quad(m\ge 1);
  \]
  \item 反过来，在固定 \(\Delta_0=c_0\) 后，\((\Delta_m)_{m\ge 0}\) 可由 \((|\alpha_m|)_{m\ge 0}\) 递推唯一恢复，具体为
  \[
  \Delta_1=c_0^2(1-|\alpha_0|^2),
  \qquad
  \Delta_{m+1}=(1-|\alpha_m|^2)\frac{\Delta_m^2}{\Delta_{m-1}}
  \quad(m\ge 1).
  \]
\end{enumerate}
\end{theorem}

\begin{proof}
式 \eqref{eq:xi-toeplitz-principal-minor-discrete-curvature} 正是推论
\ref{cor:xi-toeplitz-principal-minor-logconcavity} 的比值恒等式。
由于 \(\Delta_m>0\)，其右端属于 \((0,1]\)，故
\[
\Delta_m^2\ge \Delta_{m-1}\Delta_{m+1},
\]
且等号成立当且仅当 \(1-|\alpha_m|^2=1\)，亦即 \(|\alpha_m|=0\)。

对 \(m=0\)，由 \eqref{eq:xi-toeplitz-det-product} 在 \(m=1\) 的情形得到
\[
\Delta_1=c_0^2(1-|\alpha_0|^2),
\]
故 \(|\alpha_0|\) 由 \((\Delta_0,\Delta_1)\) 唯一确定。
对 \(m\ge 1\)，\eqref{eq:xi-toeplitz-principal-minor-discrete-curvature} 逐项给出 \(|\alpha_m|\)。
于是 \((|\alpha_m|)_{m\ge 0}\) 完全由 \((\Delta_m)_{m\ge 0}\) 唯一决定。

反向恢复则直接来自上述两条递推式：先由 \(\Delta_0=c_0\) 与 \(|\alpha_0|\) 得到 \(\Delta_1\)，再对 \(m\ge 1\) 依次使用
\[
\Delta_{m+1}=(1-|\alpha_m|^2)\frac{\Delta_m^2}{\Delta_{m-1}},
\]
即可唯一重建全部主子式。证毕。
\end{proof}

\begin{proposition}[反射谱的 \texorpdfstring{$\ell^2$}{l2} 能量与行列式离散曲率总量的等价]\label{prop:xi-toeplitz-curvature-energy-equivalence}
\leanverified{paper\_xi\_toeplitz\_curvature\_energy\_equivalence}
沿用推论 \ref{cor:xi-toeplitz-principal-minor-logconcavity} 的记号，并假设 \(\mathbf{T}_N\succ 0\) 对所有 \(N\) 成立。
令 \(\Delta_N:=\det\mathbf{T}_N>0\)，并定义离散曲率序列
\[
\mathcal{K}_N:=-\log\!\left(\frac{\Delta_{N+1}\Delta_{N-1}}{\Delta_N^{2}}\right)
=-\log(1-|\alpha_N|^2)\in[0,\infty)\qquad(N\ge 1).
\]
则成立等价关系
\[
\boxed{\;
\sum_{N\ge 1}|\alpha_N|^2<\infty
\quad\Longleftrightarrow\quad
\sum_{N\ge 1}\mathcal{K}_N<\infty.\;}
\]
\end{proposition}
\begin{proof}
记 \(x_N:=|\alpha_N|^2\in[0,1)\)，则 \(\mathcal{K}_N=-\log(1-x_N)\)。
若 \(\sum x_N<\infty\)，则 \(x_N\to 0\)，从某个 \(N_0\) 起有 \(x_N\le \tfrac12\)。
利用基本不等式
\[
x\ \le\ -\log(1-x)\ \le\ 2x\qquad(0\le x\le\tfrac12),
\]
并将有限个首项吸收到常数中，即得 \(\sum \mathcal{K}_N<\infty\)。
反向同理：若 \(\sum\mathcal{K}_N<\infty\)，则 \(\mathcal{K}_N\to 0\) 从而 \(x_N\to 0\) 并最终落入 \(x_N\le\tfrac12\) 区间；同一比较不等式给出 \(\sum x_N<\infty\)。证毕。
\end{proof}

\begin{proposition}[行列式比率极限与反射谱 \texorpdfstring{$\ell^2$}{l2} 的等价刻画]\label{prop:xi-toeplitz-determinant-ratio-limit-l2}
在命题 \ref{prop:xi-toeplitz-curvature-energy-equivalence} 的假设下，定义行列式比率
\[
r_N:=\frac{\Delta_N}{\Delta_{N-1}}\qquad(N\ge 1).
\]
则由 \eqref{eq:xi-toeplitz-det-product} 有闭式
\[
\boxed{\;
r_N=c_0\prod_{k=0}^{N-1}\bigl(1-|\alpha_k|^2\bigr).\;}
\]
因此 \((r_N)\) 单调递减且存在极限 \(r_\infty\in[0,c_0]\)。更进一步，
\[
\boxed{\;
r_\infty>0
\quad\Longleftrightarrow\quad
\sum_{k\ge 0}|\alpha_k|^2<\infty.\;}
\]
并且在 \(r_\infty>0\) 的情形下，
\[
\boxed{\;\lim_{N\to\infty}\Delta_N^{1/N}=r_\infty.\;}
\]
\end{proposition}
\leanverified{paper\_xi\_toeplitz\_determinant\_ratio\_limit\_l2}
\begin{proof}
比率公式由 \eqref{eq:xi-toeplitz-det-product} 直接计算得出。
取对数得到
\(
\log r_N=\log c_0+\sum_{k=0}^{N-1}\log(1-|\alpha_k|^2)
\)。
故 \(r_\infty>0\) 当且仅当级数 \(\sum_{k\ge 0}-\log(1-|\alpha_k|^2)\) 收敛；
由命题 \ref{prop:xi-toeplitz-curvature-energy-equivalence}，这等价于 \(\sum|\alpha_k|^2<\infty\)。
最后，\(\Delta_N=\Delta_0\prod_{n=1}^N r_n\)，取对数并用 Ces\`aro 平均得
\[
\frac{1}{N}\log\Delta_N=\frac{1}{N}\log\Delta_0+\frac{1}{N}\sum_{n=1}^N\log r_n\ \longrightarrow\ \log r_\infty,
\]
从而 \(\Delta_N^{1/N}\to r_\infty\)。证毕。
\end{proof}

\begin{remark}[与 \texorpdfstring{$\ZZ/2$}{Z/2} 扇区失败定位的相容性]\label{rem:xi-verblunsky-z2-compatible}
附录推论 \ref{cor:app-horizon-toeplitz-z2-cert-split} 给出 Toeplitz--PSD 层级在实结构下的偶/奇（\(\ZZ/2\)）证书分裂，并可定位最小失败对 \((N_\star,\epsilon)\)。
定理 \ref{thm:xi-toeplitz-det-verblunsky} 的反射系数定位提供一条更细的递归坐标：它把同一失败见证压缩到最小索引 \(m_\star\) 与单标量事件 \(|\alpha_{m_\star-1}|\ge 1\) 上。
两者本质上同源：\((N_\star,\epsilon)\) 与 \(\alpha\)-前缀都由同一截断矩数据 \((c_0,\dots,c_{N_\star})\)（从而由有限阶 \(\ell\)-喷射）可审计地生成，但后者给出更高分辨率的失败步定位。
\end{remark}

\paragraph{Toeplitz 证书的 Schur--Cohn LMI 提升}
文内对有限维通道多项式已引入 Schur--Cohn Hermitian 形式矩阵 \(H(p)\) 并以半正定性作为单位圆盘内根分布的可审计 LMI 证书（注 \ref{rem:xi-schur-cohn-sdp-interface}）。
下述结论把 Toeplitz--PSD 层级本身提升到同一 \(H(p)\) 语法中，从而统一两类证书接口。

\begin{theorem}[Carath\'eodory--Toeplitz 证书的 Schur--Cohn 提升]\label{thm:xi-toeplitz-schur-cohn-lift}
对每个 \(N\ge 0\)，由截断矩数据 \((c_0,\dots,c_N)\)（等价地由单点喷射 \((\ell_1,\dots,\ell_{N+1})\)，命题 \ref{prop:app-horizon-coeff-local}）可唯一确定一段反射系数
\(
\alpha_0,\dots,\alpha_{N-1}
\)
（当 \(\mathbf{T}_N\succ 0\) 时其落在 \(\mathbb{D}\) 内，定理 \ref{thm:xi-toeplitz-det-verblunsky}）。
令 \(p_N\) 为由该段反射系数生成的首一 Szeg\H{o} 多项式（即 Schur 递归的第 \(N\) 级预测多项式；参见 \cite{Simon2005OPUC1}），其次数为 \(N{+}1\)。
则存在可逆矩阵 \(R_N\) 使
\begin{equation}\label{eq:xi-toeplitz-schur-cohn-congruence}
\boxed{\;
R_N^\ast\,\mathbf{T}_N\,R_N
\;=\;
H(p_N)\;}
\end{equation}
其中 \(H(\cdot)\) 为注 \ref{rem:xi-schur-cohn-sdp-interface} 所定义的 Schur--Cohn Hermitian 形式矩阵。
因此
\[
\mathbf{T}_N\succeq 0
\quad\Longleftrightarrow\quad
H(p_N)\succeq 0,
\]
并且当 \(\mathbf{T}_N\) 在最小级 \(N_\star\) 失败时，\(H(p_{N_\star})\) 的负惯性指标提供同阶的代数失败见证。
\end{theorem}

\begin{proof}
该结论是 Carath\'eodory--Fej\'er 截断矩问题与 Schur 算法的标准等价：截断矩 \((c_0,\dots,c_N)\) 的可行性等价于存在 Schur 参数 \(\alpha_0,\dots,\alpha_{N-1}\) 满足 \(|\alpha_k|\le 1\)，并可由此构造第 \(N\) 级预测多项式 \(p_N\)。
式 \eqref{eq:xi-toeplitz-schur-cohn-congruence} 可视为该预测多项式的 Schur--Cohn 形式与 Toeplitz 矩形式之间的同一二次型在两组自然基下的表示变换；相应的变换矩阵 \(R_N\) 可取为 Levinson--Durbin 递推产生的三角因子（参见 \cite{DymYoung1990SchurCohnMatrixPolynomials,Simon2005OPUC1} 对 Schur--Cohn 与 Schur 参数/预测多项式之间关系的综述）。
惯性指标在合同变换下不变，从而失败见证可在 \(\mathbf{T}_N\) 与 \(H(p_N)\) 之间无损转移。证毕。
\end{proof}

\begin{theorem}[Toeplitz--PSD 最小失败级与 Schur--Cohn 根越界级的严格同一]\label{thm:xi-toeplitz-psd-minimal-failure-schur-cohn-root-escape}
\leanverified{paper\_xi\_toeplitz\_psd\_minimal\_failure\_schur\_cohn\_root\_escape}
沿用定理 \ref{thm:xi-toeplitz-schur-cohn-lift} 的记号，并假设 Toeplitz--PSD 层级并非全阶成立。
定义最小失败级
\[
N_\star:=\min\{N\ge 1:\ \mathbf{T}_N\nsucceq 0\}.
\]
则有严格等式
\[
\boxed{\;
N_\star
=
\min\{N\ge 1:\ p_N\ \text{至少有一根落在 }\CC\setminus\overline{\mathbb D}\}\; }.
\]
更一般地，对每个 \(N\ge 0\) 都有
\[
\boxed{\;
\mathbf{T}_N\succeq 0
\quad\Longleftrightarrow\quad
p_N\ \text{的全部根都落在 }\overline{\mathbb D}\; }.
\]
并且
\[
\boxed{\;
\mathbf{T}_N\succ 0
\quad\Longleftrightarrow\quad
p_N\ \text{的全部根都落在 }\mathbb D\; }.
\]
\end{theorem}

\begin{proof}
由定理 \ref{thm:xi-toeplitz-schur-cohn-lift}，存在可逆矩阵 \(R_N\) 使
\[
R_N^\ast\mathbf{T}_N R_N=H(p_N),
\]
故 \(\mathbf{T}_N\) 与 \(H(p_N)\) 具有相同惯性指数。
另一方面，Schur--Cohn 理论断言：\(H(p_N)\succeq 0\) 当且仅当 \(p_N\) 的全部根都落在闭单位圆 \(\overline{\mathbb D}\) 内，而 \(H(p_N)\succ 0\) 当且仅当全部根都落在开单位圆 \(\mathbb D\) 内。
由此立即得到两条逐级等价式。

若 Toeplitz--PSD 层级失败，则对每个 \(N<N_\star\) 都有 \(\mathbf{T}_N\succeq 0\)，从而 \(p_N\) 的全部根都落在 \(\overline{\mathbb D}\)；
而在 \(N=N_\star\) 时，\(\mathbf{T}_{N_\star}\nsucceq 0\)，故 \(p_{N_\star}\) 至少有一根落在 \(\CC\setminus\overline{\mathbb D}\)。
这正给出所述最小次数同一性。证毕。
\end{proof}

\begin{remark}[边界接触与严格正定失败]\label{rem:xi-toeplitz-strict-failure-boundary-contact}
若把“失败级”改定义为
\[
\min\{N\ge 1:\ \mathbf{T}_N\not\succ 0\},
\]
则对应的最小次数刻画的正是 \(p_N\) 首次接触单位圆边界或越出闭单位圆的层级。
因此，半正定失败控制“外逃”，而严格正定失败同时记录“边界接触或外逃”。
\end{remark}

\begin{proposition}[Toeplitz 正性在不变系数环上的区间化审计]\label{prop:xi-toeplitz-interval-audit-invariant-ring}
在定理 \ref{thm:xi-toeplitz-schur-cohn-lift} 的设定下，可取可逆下三角矩阵 \(L_N\)（例如 Levinson--Durbin 递推产生的三角因子）使
\[
\boxed{\;
\mathbf{T}_N=L_N^{\ast}\,H(p_N)\,L_N\; }.
\]
因此 \(\mathbf{T}_N\) 与 \(H(p_N)\) 具有相同惯性指数，且 Toeplitz 正性判据可完全转写为 Schur--Cohn 的半正定证书接口。
若进一步处于自对偶/实结构与不变量子环口径（系数落在 \(\QQ[u+u^{-1}]\)），则沿单位圆 \(u=e^{it}\) 可把系数视为实变量 \(x=2\cos t\in[-2,2]\) 的多项式；
从而 Schur--Cohn 证书可压缩为区间半正定约束
\[
H(x)\succeq 0\qquad(x\in[-2,2])
\]
（见注 \ref{rem:xi-schur-cohn-sdp-interface}）。
\end{proposition}
\begin{proof}
由 \eqref{eq:xi-toeplitz-schur-cohn-congruence} 得
\(
\mathbf{T}_N=R_N^{-*}H(p_N)R_N^{-1}
\)；
并且定理 \ref{thm:xi-toeplitz-schur-cohn-lift} 的证明指出 \(R_N\) 可取为三角因子，故取 \(L_N:=R_N^{-1}\) 即得所示形式。
区间化部分由注 \ref{rem:xi-schur-cohn-sdp-interface} 的降维机制直接给出。证毕。
\end{proof}

\paragraph{有限喷射对 CMV 证书与 Toeplitz 失败层级的局域闭包}
在定理 \ref{thm:xi-verblunsky-from-local-jet} 已把 Verblunsky 前缀写成有限喷射的有理函数、定理 \ref{thm:xi-toeplitz-det-verblunsky} 已把 Toeplitz 主子式与最小失败步压缩为反射系数乘积之后，有限阶喷射数据还可继续闭合为有限 CMV 证书与失败层级的严格局域律。其要点在于：长度 \(N{+}2\) 的单点喷射既已足以确定前 \(N\) 个反射参数，故也足以确定全部前 \(N\) 级 Toeplitz 主子式、相应有限 CMV 证书以及“第一失败步”在该截断范围内的位置；进一步，把这些有理依赖清分母后，整个有限喷射空间将被第一失败指标切分为两两不交的半代数层。

\begin{theorem}[有限喷射对 Verblunsky 前缀、Toeplitz 主子式与有限 CMV 证书的严格局域决定]\label{thm:xi-local-jet-determines-finite-cmv-and-toeplitz-certificate}
\leanverified{paper\_xi\_local\_jet\_determines\_finite\_cmv\_and\_toeplitz\_certificate}
沿用本段记号。对任意整数 \(N\ge 1\)，记长度 \(N{+}2\) 的单点喷射向量为
\[
J_{N+2}:=(\ell_1,\dots,\ell_{N+2}).
\]
则 \(J_{N+2}\) 唯一决定以下全部有限级对象：
\[
(\alpha_0,\dots,\alpha_{N-1}),
\qquad
(\Delta_1,\dots,\Delta_N),
\qquad
\mathcal C^{(N)},
\]
其中
\[
\Delta_m:=\det \mathbf T_m\qquad(1\le m\le N),
\]
而 \(\mathcal C^{(N)}\) 表示由 \((\alpha_0,\dots,\alpha_{N-1})\) 所生成的有限 CMV 证书。
特别地，若两组喷射数据在前 \(N{+}2\) 阶完全一致，则它们的前 \(N\) 个 Verblunsky 系数、前 \(N\) 个 Toeplitz 主子式与前 \(N\) 级 CMV 证书完全一致；并且其最小失败指数
\[
m_\star\wedge (N+1),
\qquad
m_\star:=\min\{m\ge 1:\ \mathbf T_m\nsucceq 0\},
\]
亦完全一致；此处若该集合为空，则约定 \(m_\star=\infty\)。
\end{theorem}
\begin{proof}
由命题 \ref{prop:app-horizon-coeff-local}，喷射 \(J_{N+2}\) 唯一决定 \(c_0,\dots,c_{N+1}\)。再由定理 \ref{thm:xi-verblunsky-from-local-jet}，对每个 \(0\le k\le N-1\)，\(\alpha_k\) 是 \((\ell_1,\dots,\ell_{k+2})\) 的有理函数，故 \(J_{N+2}\) 唯一决定
\[
\alpha_0,\dots,\alpha_{N-1}.
\]
有限 CMV 证书 \(\mathcal C^{(N)}\) 的矩阵元仅依赖于这段反射系数前缀，故亦被 \(J_{N+2}\) 唯一决定。

另一方面，由定理 \ref{thm:xi-toeplitz-det-verblunsky} 的乘积分解，
\[
\Delta_m
=
c_0^{m+1}\prod_{j=0}^{m-1}(1-|\alpha_j|^2)^{\,m-j}
\qquad(1\le m\le N),
\]
故 \((\Delta_1,\dots,\Delta_N)\) 也完全由 \(J_{N+2}\) 决定。

最后，由同一定理，最小失败指数 \(m_\star\) 在截断范围内恰由首个满足
\[
|\alpha_{m-1}|\ge 1
\]
的指标 \(m\) 给出；若所有 \(0\le j\le N-1\) 都满足 \(|\alpha_j|<1\)，则 \(m_\star\wedge (N+1)=N+1\)。因此 \(m_\star\wedge (N+1)\) 也仅依赖于 \((\alpha_0,\dots,\alpha_{N-1})\)，从而仅依赖于 \(J_{N+2}\)。证毕。
\end{proof}

\begin{theorem}[Toeplitz--PSD 失败指数在有限喷射空间中的半代数分层]\label{thm:xi-local-jet-toeplitz-failure-semialgebraic-stratification}
\leanverified{paper\_xi\_local\_jet\_toeplitz\_failure\_semialgebraic\_stratification}
沿用定理 \ref{thm:xi-local-jet-determines-finite-cmv-and-toeplitz-certificate} 的记号。固定 \(N\ge 1\)，定义非退化喷射空间 \(\mathcal J_{N+2}^{\mathrm{nd}}\) 为满足下列条件的实喷射向量 \(J=(\ell_1,\dots,\ell_{N+2})\) 的集合：
\begin{enumerate}
  \item \(c_0(J)>0\)；
  \item \(\alpha_0(J),\dots,\alpha_{N-1}(J)\) 的有理表达式分母均不为零。
\end{enumerate}
对 \(1\le m\le N\)，定义
\[
\mathcal P_N
:=
\{J\in \mathcal J_{N+2}^{\mathrm{nd}}:\ |\alpha_k(J)|<1\ \text{对所有 }0\le k\le N-1\},
\]
\[
\mathcal F_m
:=
\{J\in \mathcal J_{N+2}^{\mathrm{nd}}:\ |\alpha_0(J)|<1,\dots,|\alpha_{m-2}(J)|<1,\ |\alpha_{m-1}(J)|\ge 1\}.
\]
则成立：
\[
\mathcal J_{N+2}^{\mathrm{nd}}
=
\mathcal P_N\ \sqcup\ \bigsqcup_{m=1}^{N}\mathcal F_m,
\]
且这是一个两两不交的半代数分层。更具体地，记
\[
\Delta_m(J):=\det \mathbf T_m(J),
\]
则：
\begin{enumerate}
  \item \(J\in\mathcal P_N\) 当且仅当
  \[
  \Delta_1(J)>0,\dots,\Delta_N(J)>0;
  \]
  \item 对每个 \(1\le m\le N\)，\(J\in\mathcal F_m\) 当且仅当
  \[
  \Delta_1(J)>0,\dots,\Delta_{m-1}(J)>0,
  \qquad
  \Delta_m(J)\le 0;
  \]
  \item 因而函数
  \[
  \mathrm{Fail}(J):=
  \min\{m\ge 1:\ \mathbf T_m(J)\nsucceq 0\}\wedge (N+1)
  \]
  （若该集合为空则约定其最小值为 \(\infty\)）是 \(\mathcal J_{N+2}^{\mathrm{nd}}\) 上的有限值半代数函数。
\end{enumerate}
\end{theorem}
\begin{proof}
由定理 \ref{thm:xi-verblunsky-from-local-jet}，每个 \(\alpha_k(J)\) 都是喷射坐标的有理函数；在 \(\mathcal J_{N+2}^{\mathrm{nd}}\) 上分母不为零，因此条件 \(|\alpha_k(J)|<1\) 与 \(|\alpha_k(J)|\ge 1\) 清分母后都可写成有限个多项式不等式。故 \(\mathcal P_N\) 与各 \(\mathcal F_m\) 均为半代数集。

它们两两不交由定义立即可见。再证覆盖性：任取 \(J\in \mathcal J_{N+2}^{\mathrm{nd}}\)。若 \(|\alpha_k(J)|<1\) 对所有 \(0\le k\le N-1\) 都成立，则 \(J\in \mathcal P_N\)；否则存在最小 \(m\in\{1,\dots,N\}\) 使 \(|\alpha_{m-1}(J)|\ge 1\)，于是 \(J\in \mathcal F_m\)。故得到所述不交并分解。

对第 \(1\) 条，若 \(J\in \mathcal P_N\)，则由定理 \ref{thm:xi-toeplitz-det-verblunsky} 的乘积分解
\[
\Delta_m(J)
=
c_0(J)^{m+1}\prod_{j=0}^{m-1}(1-|\alpha_j(J)|^2)^{\,m-j}
\]
可知 \(\Delta_m(J)>0\) 对所有 \(1\le m\le N\) 成立。反之，若 \(\Delta_1(J),\dots,\Delta_N(J)\) 全部为正，而某个最小 \(m\le N\) 满足 \(|\alpha_{m-1}(J)|\ge 1\)，则由定理 \ref{thm:xi-toeplitz-det-verblunsky} 的最小失败步定位必有 \(\Delta_m(J)\le 0\)，矛盾。故必有 \(J\in \mathcal P_N\)。

对第 \(2\) 条，若 \(J\in \mathcal F_m\)，则前 \(m-1\) 个反射系数均落在单位圆盘内，故由同一乘积分解可得
\[
\Delta_1(J),\dots,\Delta_{m-1}(J)>0.
\]
又因 \(|\alpha_{m-1}(J)|\ge 1\)，故
\[
1-|\alpha_{m-1}(J)|^2\le 0,
\]
从而 \(\Delta_m(J)\le 0\)。反之，若
\[
\Delta_1(J),\dots,\Delta_{m-1}(J)>0,
\qquad
\Delta_m(J)\le 0,
\]
则由第 \(1\) 条应用于 \(m-1\) 级可知
\[
|\alpha_0(J)|<1,\dots,|\alpha_{m-2}(J)|<1.
\]
若再有 \(|\alpha_{m-1}(J)|<1\)，则乘积分解会推出 \(\Delta_m(J)>0\)，矛盾。故必有 \(|\alpha_{m-1}(J)|\ge 1\)，即 \(J\in \mathcal F_m\)。

第 \(3\) 条只是把上述不交并分解重新编码为失败指数函数的纤维：\(\mathrm{Fail}^{-1}(N+1)=\mathcal P_N\)，而对每个 \(1\le m\le N\) 有 \(\mathrm{Fail}^{-1}(m)=\mathcal F_m\)。因此 \(\mathrm{Fail}\) 是有限值半代数函数。证毕。
\end{proof}

\begin{theorem}[Toeplitz--PSD 层级的遗传性、共终稀疏核验与全称量词消去]\label{thm:xi-toeplitz-psd-cofinal-sparsification-hereditary}
\leanverified{paper\_xi\_toeplitz\_psd\_cofinal\_sparsification\_hereditary}
沿用本段 Toeplitz 主子矩阵
\[
\mathbf T_N=(c_{i-j})_{0\le i,j\le N}
\]
的记号，并定义谓词
\[
P_{\mathrm{PSD}}(N):\qquad \mathbf T_N\succeq 0.
\]
则：
\begin{enumerate}
  \item \(P_{\mathrm{PSD}}(N)\) 是遗传谓词，即对任意 \(0\le N\le N'\)，若 \(\mathbf T_{N'}\succeq 0\)，则 \(\mathbf T_N\succeq 0\)；
  \item 对任意严格递增且共终的整数序列 \((N_j)_{j\ge 1}\)（即 \(N_j\to\infty\)），有严格等价
  \[
  \forall N\ge 0,\ \mathbf T_N\succeq 0
  \quad\Longleftrightarrow\quad
  \forall j\ge 1,\ \mathbf T_{N_j}\succeq 0;
  \]
  \item 更一般地，若 \(\Phi(N)\) 是任一遗传谓词，则同样有
  \[
  \forall N\ge 0,\ \Phi(N)
  \quad\Longleftrightarrow\quad
  \forall j\ge 1,\ \Phi(N_j).
  \]
  特别地，由有限个遗传谓词经有限次合取与析取所构成、且结果仍为遗传的层级断言，都允许把“对所有深度 \(N\)”的全称量词替换为对任意共终稀疏子序列的全称量词。
\end{enumerate}
\end{theorem}

\begin{proof}
第 \(1\) 条是 Hermitian 正半定性对主子矩阵的稳定性：\(\mathbf T_N\) 本就是 \(\mathbf T_{N'}\) 的左上角主子矩阵，故若 \(\mathbf T_{N'}\succeq 0\)，则 \(\mathbf T_N\succeq 0\)。

对第 \(2\) 条，正向 implication 显然。反向设 \(\mathbf T_{N_j}\succeq 0\) 对所有 \(j\) 成立。任取 \(N\ge 0\)。由于 \((N_j)\) 共终，存在 \(j\) 使 \(N\le N_j\)。由第 \(1\) 条的遗传性，\(\mathbf T_{N_j}\succeq 0\) 推出 \(\mathbf T_N\succeq 0\)。故 \(\mathbf T_N\succeq 0\) 对所有 \(N\) 都成立。

第 \(3\) 条与第 \(2\) 条完全同理：若 \(\Phi\) 本身遗传，则由共终性可对任意固定 \(N\) 选取 \(j\) 使 \(N\le N_j\)，再由
\[
\Phi(N_j)\Longrightarrow \Phi(N)
\]
得到反向 implication。特别地，有限个遗传谓词的合取与析取仍为遗传，故同一量词消去结论继续成立。证毕。
\end{proof}


\begin{remark}[与有限型 CMV/Schur 接口的对齐]\label{rem:xi-verblunsky-cmv-align-finite}
命题 \ref{prop:xi-hilbert-polya-cmv} 在有限窗/有限维口径下把“单位圆根条件”转写为 CMV-like 酉谱实现，并以 \(|\alpha_k|<1\) 作为可审计接口。
本段的构造给出其无穷维/极限版：当 RH 成立时，\(\mathscr{C}_\zeta\) 的反射谱 \((\alpha_n)\subset\mathbb{D}\) 生成一个规范 CMV 算子，其谱测度即为视界正测度 \(\mu_\zeta\)（附录定理 \ref{thm:app-horizon-spectral-measure-atomic}）。
因此，“Toeplitz--PSD 层级的有限级证书”可被等价理解为“CMV 反射谱的前缀约束”，而其失败可由最小越界反射系数 \(|\alpha_{m_\star-1}|\ge 1\) 精确定位（定理 \ref{thm:xi-toeplitz-det-verblunsky}）。
\end{remark}

\endinput
